\subsection{波莱尔截面}

定理 4. 设 X 是波兰空间，R 是 X 上的一个等价关系，使得模 R 的等价类在 X 中闭，且 X 中每个闭集（关于 R ）的饱和集是波莱尔集。则存在 X 中的一个波莱尔集，它与每个等价类恰交于一点。

取 X 上与它的拓扑相容且使 X 完备的一个度量。由第 5 号的引理 3 ，存在 X 的一个筛分，由一个筛 \(\mathbf{C} = (\mathbf{C}_{n}, p_{n})\) 与一个映射序列 \((\varphi_{n})\) 定义。对每个 \(c \in C_{n}\) ，令 \(g_{n}(c)\) 为闭集 \(\varphi_{n}(c)\) 关于 R 的饱和集；由假设，\(g_{n}(c)\) 是 X 中的波莱尔集。

由于每个集 \(\mathbf{C}_n\) 可数，我们可以把每个 \(\mathbf{C}_n\) 线性排序，使得小于一个给定元素的元素之集有限。对每个 \(c \in \mathbf{C}_n\) ，我们定义一个集 \(h_n(c)\) ，用对 \(n\) 的归纳，如下。首先，对 \(c \in \mathbf{C}_0\) ，\(h_0(c)\) 是 \(\varphi_0(c)\) 与诸集 \(\mathbf{X} - g_0(c')\) （其中 \(c' \in \mathbf{C}_0\) 且 \(c' < c\) ）的交。对 \(c \in \mathbf{C}_{n+1}\) ，\(h_{n+1}(c)\) 是 \(\varphi_{n+1}(c)\) 、\(h_n(\mathbf{P}_n(c))\) 与诸集 \(\mathbf{X} - g_{n+1}(c')\) （对 \(c' \in \mathbf{C}_{n+1}\) ，\(p_n(c') = p_n(c)\) 且 \(c' < c\) ）的交。诸 \(h_n(c)\) 显然是波莱尔集。

我们将证明下述断言：对每个整数 \(n \geqslant 0\) 与每个等价类 \(\mathbf{H} \mod \mathbb{R}\) ，存在唯一的元素 \(c \in \mathbf{C}_n\) 使 \(h_n(c)\) 与 \(\mathbf{H}\) 相交，且我们有

\[
h (c _ {n}) \cap \mathrm{H} = \varphi_ {n} (c) \cap \mathrm{H},
\]

它因而是闭集。对 \(n = 0\) ，考虑最小的元素 \(c \in \mathbf{C}_0\) （使 \(\varphi_0(c)\) 与 \(\mathbf{H}\) 相交者）；则 \(\varphi_0(c) \cap \mathbf{H}\) 不交任何集 \(g_0(c')\) （对满足 \(c' \in \mathbf{C}_0\) 且 \(c' < c\) 的 c' ）；故它含于 \(h_0(c) \cap \mathbf{H}\) ，从而等于此集；此外，我们有 \(\mathbf{H} \subset g_0(c)\) ，因此 \(h_0(c') \cap \mathbf{H}\) 为空（对 \(c' \in \mathbf{C}_0\) 且 \(c' > c\) ）；这样断言对 \(n = 0\) 得证。我们对 \(n\) 继续用归纳法：若存在 \(c \in \mathbf{C}_{n+1}\) 使 \(h_{n+1}(c)\) 与 \(\mathbf{H}\) 相交，则由关系 \(h_{n+1}(c) \subset h_n(p_n(c))\) 与归纳假设，\(p_n(c)\) 是唯一的元素 \(d \in \mathbf{C}_n\) （使 \(h_n(d)\) 与 \(\mathbf{H}\) 相交者）。注意 \(h_n(d)\) 含于 \(\varphi_n(d)\) ，由筛分的定义，它含于诸集 \(\varphi_n(c)\) （对 \(c \in p_n(d)\) ）的并；因此存在最小元素 \(c \in p_n(d)\) 使 \(\varphi_{n+1}(c)\) 与 \(\mathbf{H}\) 相交。因此我们有

\[
\varphi_ {n + 1} (c) \cap \mathrm{H} \subset \varphi_ {n} (d) \cap \mathrm{H} = h _ {n} (d) \cap \mathrm{H}
\]

（依归纳假设）。从而

\[
\varphi_ {n + 1} (c) \cap \mathrm{H} \subset \varphi_ {n + 1} (c) \cap h _ {n} (d),
\]

且由于依定义 \(\varphi_{n+1}(c) \cap \mathbf{H}\) 不交任何集 \(g_{n+1}(c')\) （对满足 \(c' \in p_n(d)\) 且 \(c' < c\) 的 c' ），由 \(h_{n+1}(c)\) 的定义即得 \(\varphi_{n+1}(c) \cap \mathbf{H} = h_{n+1}(c) \cap \mathbf{H}\) 。此外，我们有 \(\mathbf{H} \subset g_{n+1}(c)\) ，因此若 \(c' \in p_n(d)\) 满足 \(c' > c\) ，则 \(h_{n+1}(c') \cap \mathbf{H}\) 为空。这样断言对一切 \(n\) 得证。

对每个整数 \(n\) ，令 \(\mathbf{S}_n\) 为诸集 \(h_n(c)\) （其中 \(c\) 跑遍 \(\mathbf{C}_n\) ）的并。集 \(\mathbf{S}_n\) 是波莱尔集，且我们有 \(\mathbf{S}_{n+1} \subset \mathbf{S}_n\) 。

令 S 为诸集 \(S_{n}\) 的交，它也是 X 中的波莱尔集；我们将证明 S 与每个等价类 H 模 R 恰交于一点。对每个 n ，令 \(c_{n}(H)\) 为唯一的元素 \(c \in C_{n}\) （使 \(h_{n}(c)\) 与 H 相交者）；则 \(\mathbf{S}_{n} \cap \mathbf{H} = \varphi_{n}(c_{n}(\mathbf{H})) \cap \mathbf{H}\) ，而 \(S \cap H\) 是诸集 \(\varphi_{n}(c_{n}(\mathbf{H})) \cap \mathbf{H}\) 的交。由于序列 \((c_{n}(\mathbf{H}))\) 属于 L(C)，闭集的递减序列 \(\varphi_{n}(c_{n}(\mathbf{H}))\) —— 其直径趋于 0 —— 收敛于一点 \(x \in X\) ，因为 X 完备。诸闭集 \(\varphi_{n}(c_{n}(\mathbf{H})) \cap \mathbf{H}\) 的交因此由单独一点 x 组成，定理 4 的证明完成。

注. 特别地，闭等价关系 R 满足定理 4 的假设。当 X 是紧可度量化空间时，定理 4 因此适用于任意豪斯多夫等价关系 R ，因为紧空间上的豪斯多夫等价关系是闭的（第 I 章，§ 10，第 4 号，命题 8）。

